\section{导读：为什么 Lovart 是垂直 Agent 样本}

本节先定位这期。Lovart 是面向设计师的垂直 Agent，而不是 Manus 那类通用 Agent。它的价值不在“什么都能做”，而是在设计工作流里理解 brief、素材、审美、版本迭代和交付。陈冕的访谈特别之处，是它同时呈现了 AI 应用创业的产品逻辑和创始人的情绪状态：补贴战、下架、现金见底、融不到资之后，终于出现全球传播和产品认可。

\begin{figure}[H]
\centering
\includegraphics[width=0.94\textwidth]{lovart-positioning.png}
\caption{Lovart：垂直 Agent。Manus 是通用 Agent，Lovart 面向设计师和创意工作流。自制概念图，依据视频描述和 00:54:56--01:10:30 对谈内容整理。}
\end{figure}

\begin{knowledgebox}{读图：通用 Agent 与垂直 Agent 的差别}
通用 Agent 追求任务覆盖面，垂直 Agent 追求在一个专业场景里把上下文、工具、审美和交付打深。Lovart 的判断是：设计师工作流足够复杂，值得专门做一个 Agent。
\end{knowledgebox}

\begin{importantbox}{本期核心命题}
AI 应用创业的关键不是“接一个更强模型”，而是在模型能力窗口里找到高价值工作流，并把 Magic Moment 转化成留存、付费和专业心智。
\end{importantbox}

\begin{knowledgebox}{本期阅读路线}
\begin{tabularx}{\textwidth}{p{0.22\textwidth}p{0.34\textwidth}X}
\toprule
路线 & 主要问题 & 为什么重要 \\
\midrule
产品路线 & Lovart 为什么选择设计师和垂直 Agent？ & 决定它是否能避开通用 Agent 的正面竞争。 \\
工作流路线 & 从 brief 到交付到底有哪些环节？ & 决定产品是否进入真实生产。 \\
创业路线 & 补贴战、下架、4000 块现金如何影响判断？ & 决定团队是否能活到窗口打开。 \\
PMF 路线 & “好爽”如何转成留存和收入？ & 决定情绪峰值能否变成生意。 \\
壁垒路线 & 模型红利之后应用公司还剩什么？ & 决定长期竞争位置。 \\
\bottomrule
\end{tabularx}
\end{knowledgebox}

\subsection{本章小结}

Lovart 特辑是 AI 应用创业切片。它既讲垂直 Agent 产品定位，也讲创业低谷后的爆发情绪，是理解 2025 年 Agent 应用窗口的一个样本。

